% LuaLaTeX document — compile with: lualatex inter_voxel_propagator_plan.tex
\documentclass[11pt,a4paper]{article}

\usepackage{fontspec}
\setmainfont{Latin Modern Roman}
\setmonofont{Latin Modern Mono}

\usepackage{amsmath,amssymb,amsthm,bm}
\usepackage{booktabs}
\usepackage{geometry}
\geometry{margin=2.5cm}
\usepackage{hyperref}
\hypersetup{colorlinks=true,linkcolor=blue!60!black,citecolor=green!50!black,urlcolor=blue!70!black}
\usepackage{enumitem}
\usepackage{xcolor}
\usepackage{mathtools}

\newcommand{\bR}{\bm{R}}
\newcommand{\bk}{\bm{k}}
\newcommand{\bw}{\bm{w}}
\newcommand{\bs}{\bm{s}}
\newcommand{\kh}{\hat{\bm{k}}}
\newcommand{\order}{\mathcal{O}}
\newcommand{\D}{\Delta}

\theoremstyle{plain}
\newtheorem{theorem}{Theorem}
\newtheorem{proposition}[theorem]{Proposition}
\theoremstyle{remark}
\newtheorem{remark}[theorem]{Remark}

\title{Analytical Inter-Voxel Strain Propagator\\
\large Implementation Plan for the Foldy--Lax Equation}
\author{Research Notes}
\date{April 2026}

\begin{document}
\maketitle
\tableofcontents
\vspace{1em}

%======================================================================
\section{The problem: UV divergence of the strain propagator}
\label{sec:problem}
%======================================================================

The Foldy--Lax multiple-scattering equation
$\mathbf{T} = \mathbf{T}_0(\mathbf{I} - \mathbf{G}_0\mathbf{T}_0)^{-1}$
requires the \emph{inter-site strain propagator}---the volume-averaged
Green's tensor second derivative between two cubic voxels separated by~$\bR$:
\begin{equation}
  P_{ijkl}(\bR,\omega) \;=\;
  \int_A\!\int_B
  \frac{\partial^2 G_{ik}(\bs-\bs'+\bR,\,\omega)}%
       {\partial x_j\,\partial x_l}\,d^3s\,d^3s'.
  \label{eq:Pijkl-def}
\end{equation}
In Fourier space this becomes
\begin{equation}
  P_{ijkl}(\bR,\omega) \;=\;
  \int\!\frac{d^3k}{(2\pi)^3}\;
  \bigl|\tilde{f}_0(\bk)\bigr|^2\,
  \tilde{\Gamma}_{ijkl}(\bk,\omega)\,
  e^{i\bk\cdot\bR},
  \label{eq:Pijkl-fourier}
\end{equation}
where $\tilde{f}_0(\bk) = a^3\,\mathrm{sinc}(k_1 h)\,
\mathrm{sinc}(k_2 h)\,\mathrm{sinc}(k_3 h)/(2\sqrt{\pi})$
is the cube form factor ($h=a/2$) and
$\tilde{\Gamma}_{ijkl}$ is the strain kernel derived from the
Green's tensor.

\paragraph{The divergence.}
For the static Kelvin Green's tensor, the strain kernel
\begin{equation}
  \tilde{\Gamma}^\infty_{ijkl}(\kh) \;=\;
  -\frac{1}{2\mu}\bigl[
    \hat{k}_j\hat{k}_l\,\delta_{ik}
  + \hat{k}_i\hat{k}_l\,\delta_{jk}
  - 2\eta\,\hat{k}_i\hat{k}_j\hat{k}_k\hat{k}_l
  \bigr]
  \label{eq:Gamma-inf}
\end{equation}
depends only on the \emph{direction}~$\kh$, not on $|\bk|$.
It is $\order(1)$ as $|\bk|\to\infty$, which---despite the
$\mathrm{sinc}^2$ decay of $|\tilde{f}_0|^2 \sim 1/k^2$ per
axis---leaves the 3D integral \eqref{eq:Pijkl-fourier}
\textbf{UV-divergent}.  Confirmed numerically (WS1): the integral
does not converge as $K_{\max}\to\infty$.


%======================================================================
\section{Static piece: potential derivatives and delta-function collapse}
\label{sec:static}
%======================================================================

\subsection{Reduction to convergent potentials}

The key observation (WS2) is that $\kh_j\kh_l = k_j k_l / k^2$
and $\kh_i\kh_j\kh_k\kh_l = k_i k_j k_k k_l / k^4$ in the
Fourier integral can be traded for
\emph{real-space derivatives} of convergent potentials.  Define:
\begin{align}
  A_{jl}(\bR) &\;=\;
    \int\!\frac{d^3k}{(2\pi)^3}\;
    |\tilde{f}_0|^2\,\hat{k}_j\hat{k}_l\,e^{i\bk\cdot\bR}
  \;=\; -\frac{\partial^2}{\partial R_j\,\partial R_l}
  \biggl[\frac{\Phi_{00}(\bR)}{4\pi}\biggr],
  \label{eq:Ajl}
  \\[4pt]
  B_{ijkl}(\bR) &\;=\;
    \int\!\frac{d^3k}{(2\pi)^3}\;
    |\tilde{f}_0|^2\,\hat{k}_i\hat{k}_j\hat{k}_k\hat{k}_l\,
    e^{i\bk\cdot\bR}
  \;=\; -\frac{\partial^4}{\partial R_i\,\partial R_j\,
    \partial R_k\,\partial R_l}
  \biggl[\frac{\Psi_{00}(\bR)}{8\pi}\biggr],
  \label{eq:Bijkl}
\end{align}
where the \textbf{Newton potential} and \textbf{biharmonic potential}
are the tent-weighted integrals
\begin{align}
  \Phi_{00}(\bR) &\;=\;
  \int M_{00}(w_1-R_1)\,M_{00}(w_2-R_2)\,M_{00}(w_3-R_3)\;
  \frac{1}{|\bw|}\;d^3w,
  \label{eq:Phi00}
  \\[2pt]
  \Psi_{00}(\bR) &\;=\;
  \int M_{00}(w_1-R_1)\,M_{00}(w_2-R_2)\,M_{00}(w_3-R_3)\;
  |\bw|\;d^3w.
  \label{eq:Psi00}
\end{align}
Here $M_{00}(d) = \max(0, 1-|d|)$ is the tent function (cross-correlation
of two box functions of width~$a$).  Both integrals converge absolutely.

The static propagator assembles as
\begin{equation}
  \boxed{
  P^{(0)}_{ijkl}(\bR) \;=\;
  -\frac{1}{2\mu}\bigl[
    \delta_{ik}\,A_{jl}
  + \delta_{jk}\,A_{il}
  - 2\eta\,B_{ijkl}
  \bigr]
  }
  \label{eq:Pstatic}
\end{equation}
with $\eta = 1/[2(1-\nu)]$.

\subsection{Delta-function collapse}

The derivatives in \eqref{eq:Ajl}--\eqref{eq:Bijkl} act on the
tent functions, not on the kernel.  The second derivative of the
tent is a sum of delta functions:
\begin{equation}
  M_{00}''(d) \;=\; -2\,\delta(d) + \delta(d-1) + \delta(d+1).
  \label{eq:tent-delta}
\end{equation}
Each delta \textbf{collapses one integration dimension} analytically.

\paragraph{Single delta collapse (for $A_{jl}$).}
Differentiating $\Phi_{00}$ twice with respect to~$R_j$ produces
deltas in the $w_j$~integration, leaving a \textbf{2D tent-weighted
master integral}:
\begin{equation}
  I^A_2(p,q,c) \;=\;
  \int_0^1\!\int_0^1
  \frac{u^p\,v^q}{\sqrt{c^2 + u^2 + v^2}}\;du\,dv.
  \label{eq:I2A}
\end{equation}
For the biharmonic kernel:
\begin{equation}
  I^B_2(p,q,c) \;=\;
  \int_0^1\!\int_0^1
  \frac{c^2\,u^p\,v^q}{(c^2 + u^2 + v^2)^{3/2}}\;du\,dv,
  \qquad
  I^\Psi_2(p,q,c) \;=\;
  \int_0^1\!\int_0^1
  u^p\,v^q\,\sqrt{c^2 + u^2 + v^2}\;du\,dv.
  \label{eq:I2BPsi}
\end{equation}

\paragraph{Double delta collapse (for $B_{ijkl}$).}
When two derivatives act on different tent axes, both dimensions
collapse, leaving a \textbf{1D elementary integral}:
\begin{equation}
  I_1(C) \;=\;
  \int_0^1 (1-t)\,\sqrt{C + t^2}\;dt
  \;=\; \frac{1}{2}\!\left[\sqrt{C+1} + C\,
  \mathrm{arcsinh}\!\left(\frac{1}{\sqrt{C}}\right)\right]
  - \frac{1}{3}\!\left[(C+1)^{3/2} - C^{3/2}\right].
  \label{eq:I1exact}
\end{equation}

\paragraph{Sequential integration.}
All 2D masters are computed analytically by Mathematica via the proven
sequential integration pattern:
\begin{equation}
  \text{inner:}\;
  \int_0^1 \frac{u^p\,v^q}{\sqrt{c^2+u^2+v^2}}\,dv
  \;\;\xrightarrow{\text{Assumptions}}\;\;
  \text{outer:}\;
  \int_0^1 (\cdots)\,du
  \;\;\xrightarrow{\text{Simplify}}\;\;
  \text{exact closed form.}
\end{equation}
Flags: \texttt{Assumptions -> u > 0 \&\& c >= 0},
\texttt{GenerateConditions -> False} at every step.

\subsection{Status: face-adjacent propagator complete}

\texttt{InterVoxelPropagator.wl} computes the face-adjacent
propagator $P^{(0)}_{ijkl}(\bR)$ for $\bR = (a,0,0)$ with $C_{4v}$
symmetry.  All independent components ($P_{1111}$, $P_{1122}$,
$P_{2222}$, $P_{2233}$, $P_{1212}$, $P_{2323}$) obtained as exact
closed forms, validated via:
\begin{itemize}[nosep]
  \item Laplacian identity:
    $B_{1111} + B_{1122} + B_{1133} = A_{11}$
    (from $\nabla^2|\bw| = 2/|\bw|$).
  \item Finite-difference cross-check against WS2 (6-digit agreement).
\end{itemize}


%======================================================================
\section{Dynamic corrections: analytical power series in $(ka)^2$}
\label{sec:dynamic}
%======================================================================

\subsection{Expansion of the dynamic Green's tensor}

The full frequency-dependent Green's tensor in Fourier space is
\begin{equation}
  \tilde{G}_{ij}(\bk,\omega) \;=\;
  \frac{\hat{k}_i\hat{k}_j}%
       {\rho(\omega^2 - c_p^2 k^2)}
  \;+\;
  \frac{\delta_{ij} - \hat{k}_i\hat{k}_j}%
       {\rho(\omega^2 - c_s^2 k^2)}.
  \label{eq:Gij-dynamic}
\end{equation}
Expanding for large~$k$ (i.e.\ $\omega^2/(c^2 k^2) \ll 1$):
\begin{equation}
  \frac{1}{\omega^2 - c^2 k^2}
  \;=\;
  -\frac{1}{c^2 k^2}
  \sum_{n=0}^{\infty}
  \left(\frac{\omega^2}{c^2 k^2}\right)^{\!n}
  \;=\;
  -\sum_{n=0}^{\infty}
  \frac{\omega^{2n}}{c^{2n+2}\,k^{2n+2}}.
  \label{eq:geometric-series}
\end{equation}
The strain kernel $\tilde{\Gamma}_{ijkl} = -\tfrac{1}{2}(k_j k_l
\tilde{G}_{ik} + k_i k_l \tilde{G}_{jk})$ therefore has the expansion
\begin{equation}
  \boxed{
  \tilde{\Gamma}_{ijkl}(\bk,\omega)
  \;=\;
  \tilde{\Gamma}^\infty_{ijkl}(\kh)
  \;+\;
  \sum_{n=1}^{\infty}
  \omega^{2n}\,
  \frac{\tilde{\Gamma}^{(n)}_{ijkl}(\kh)}{k^{2n}}
  }
  \label{eq:Gamma-expansion}
\end{equation}
where each $\tilde{\Gamma}^{(n)}(\kh)$ is a known degree-4 tensor
polynomial in~$\kh$, with coefficients depending on $c_p$, $c_s$,
$\rho$.  The $n=0$ term is exactly the static Eshelby
kernel~\eqref{eq:Gamma-inf}.

\subsection{The potential hierarchy}

Substituting \eqref{eq:Gamma-expansion} into the Fourier
integral~\eqref{eq:Pijkl-fourier} gives
\begin{equation}
  P_{ijkl}(\bR,\omega)
  \;=\;
  \sum_{n=0}^{\infty}
  \omega^{2n}\,P^{(n)}_{ijkl}(\bR)
  \label{eq:P-series}
\end{equation}
where
\begin{equation}
  P^{(n)}_{ijkl}(\bR)
  \;=\;
  \int\!\frac{d^3k}{(2\pi)^3}\;
  |\tilde{f}_0|^2\;
  \frac{\tilde{\Gamma}^{(n)}_{ijkl}(\kh)}{k^{2n}}\;
  e^{i\bk\cdot\bR}.
  \label{eq:Pn-fourier}
\end{equation}
Since $\tilde{\Gamma}^{(n)}(\kh)$ has the same angular structure as
$\tilde{\Gamma}^\infty(\kh)$ (degree-4 in~$\kh$), each $P^{(n)}$
reduces to derivatives of \textbf{higher-order potentials}:
\begin{align}
  W_m(\bR) &\;=\;
  \int M_{00}(w_1-R_1)\,M_{00}(w_2-R_2)\,M_{00}(w_3-R_3)\;
  |\bw|^{2m-1}\;d^3w,
  \label{eq:Wm-def}
\end{align}
using the Fourier-transform hierarchy
\begin{equation}
  \int\!\frac{d^3k}{(2\pi)^3}\;
  \frac{|\tilde{f}_0|^2}{k^{2m}}\;e^{i\bk\cdot\bR}
  \;=\; c_m\,W_m(\bR)
  \label{eq:FT-hierarchy}
\end{equation}
with known normalisation constants $c_m$.
The first few potentials are:
\begin{center}
\begin{tabular}{cccl}
\toprule
$m$ & Kernel & Potential & Role \\
\midrule
$1$ & $1/|\bw|$  & $\Phi_{00}$ (Newton)     & Static: $A_{jl}$ \\
$2$ & $|\bw|$    & $\Psi_{00}$ (biharmonic)  & Static: $B_{ijkl}$; Dynamic: $P^{(1)}$ \\
$3$ & $|\bw|^3$  & $\mathrm{X}_{00}$        & Dynamic: $P^{(1)}, P^{(2)}$ \\
$4$ & $|\bw|^5$  & $\Omega_{00}$             & Dynamic: $P^{(2)}, P^{(3)}$ \\
\bottomrule
\end{tabular}
\end{center}

\begin{remark}[Smoothness improves with order]
The Newton potential ($m=1$) has a $1/|\bw|$ singularity at the
origin---this is the \emph{hardest} integral.  All higher potentials
($m\geq 2$) have smooth, bounded kernels $|\bw|^{2m-1}$.  Their
tent-weighted integrals are strictly easier for Mathematica to
evaluate in closed form.
\end{remark}

\subsection{Explicit structure of the first dynamic correction}

The order $n=1$ contribution to the strain kernel is
\begin{equation}
  \tilde{\Gamma}^{(1)}_{ijkl}(\kh)
  \;=\; \frac{1}{2}\biggl[
    \hat{k}_j\hat{k}_l\,\delta_{ik}
    \!\left(\frac{1}{\rho c_p^4}\hat{k}_j^2
    + \frac{1}{\rho c_s^4}(\delta_{jj'}-\hat{k}_j^2)\right)
    + \cdots
  \biggr],
\end{equation}
which has the same tensorial form as~\eqref{eq:Gamma-inf} but with
\emph{different} coefficients involving $c_p^{-4}$ and $c_s^{-4}$
instead of $c_p^{-2}$ and $c_s^{-2}$.  The real-space reduction
follows the same pattern:
\begin{equation}
  P^{(1)}_{ijkl}(\bR)
  \;=\;
  -\frac{1}{2}\biggl[
    \delta_{ik}\,\frac{\partial^2}{\partial R_j\partial R_l}
    \!\left(\frac{\Psi_{00}}{c_1}\right)
  + \delta_{jk}\,\frac{\partial^2}{\partial R_i\partial R_l}
    \!\left(\frac{\Psi_{00}}{c_1}\right)
  - 2\eta'\,\frac{\partial^4}{\partial R_i\partial R_j
    \partial R_k\partial R_l}
    \!\left(\frac{\mathrm{X}_{00}}{c_2}\right)
  \biggr]
  \label{eq:P1-assembly}
\end{equation}
where $c_1$, $c_2$, $\eta'$ are known functions of $c_p, c_s, \rho, \mu$.

\paragraph{Delta-function collapse applies identically.}
The tent derivatives produce the same delta sums~\eqref{eq:tent-delta}.
For $\partial^2\Psi_{00}/\partial R_j^2$, the collapse gives
2D~integrals of type $I_2^\Psi(p,q,c)$, which are \emph{already
computed} in the face-adjacent script.  For $\partial^4
\mathrm{X}_{00}/\partial R_i\cdots\partial R_l$, the collapse gives
2D~integrals of the form
\begin{equation}
  I_2^{\mathrm{X}}(p,q,c)
  \;=\;
  \int_0^1\!\int_0^1
  u^p\,v^q\,(c^2+u^2+v^2)^{3/2}\;du\,dv,
  \label{eq:I2X}
\end{equation}
which is \emph{polynomial~$\times$ smooth} and trivially yields exact
closed forms.  Similarly for double-delta collapse to 1D.

\subsection{Equivalent real-space perspective}

The Fourier series \eqref{eq:Gamma-expansion} is equivalent to
expanding the real-space Green's function kernel $e^{ik\rho}/\rho$
(where $\rho = |\bw|$) inside the tent-weighted integral:
\begin{equation}
  \frac{e^{ik\rho}}{\rho}
  \;=\;
  \frac{1}{\rho}
  + ik
  + \frac{(ik)^2}{2!}\,\rho
  + \frac{(ik)^3}{3!}\,\rho^2
  + \frac{(ik)^4}{4!}\,\rho^3
  + \cdots
  \label{eq:exp-expansion}
\end{equation}
Each term generates an integral with kernel $\rho^{n-1}$:
\begin{itemize}[nosep]
  \item $n=0$: $1/\rho$ --- the Newton potential master integrals
    $I^A_2$ (singular, already computed).
  \item $n=1$: constant --- trivially polynomial, exact rational result.
  \item $n=2$: $\rho$ --- the biharmonic master integrals
    $I^\Psi_2$ (smooth, already computed via the identity
    $I^\Psi_2 = c^2 I^A_2 + I^A_2[\text{shifted degrees}]$).
  \item $n=3$: $\rho^2$ --- polynomial, exact rational.
  \item $n=4$: $\rho^3$ --- smooth, new but easy.
  \item $n \geq 5$: increasingly smooth.
\end{itemize}
Odd-$n$ terms (constant, $\rho^2$, $\rho^4$, \ldots) are
\textbf{purely polynomial} and integrate to exact rationals.  Even-$n$
terms involve the algebraic kernels $\rho^{2m+1}$ and yield closed
forms in $\mathbb{Q}(\sqrt{2},\sqrt{3},\pi,\ldots)$---the same
number field as the existing master integrals.

\subsection{Convergence}

The series~\eqref{eq:P-series} converges for
\begin{equation}
  \frac{\omega^2}{c_s^2 k_{\mathrm{eff}}^2} < 1,
  \qquad\text{i.e.}\qquad
  ka \lesssim \pi.
\end{equation}
Here $k_{\mathrm{eff}} \sim \pi/a$ is the effective wavenumber sampled
by the sinc form factor.  In the Rayleigh regime ($ka < 0.3$) only
$P^{(0)}$ and $P^{(1)}$ are needed.  In the transition regime
($0.3 < ka < 1$), the first three or four terms suffice.
For $ka > 1$, the resonance T-matrix already subdivides the cube into
$n^3$~sub-cells with effective $ka' = ka/n < 1$.


%======================================================================
\section{Three neighbour types and their symmetry}
\label{sec:neighbours}
%======================================================================

On a cubic lattice, the 26 nearest neighbours fall into three
orbits under the octahedral group~$O_h$:
\begin{center}
\begin{tabular}{lcccc}
\toprule
Type & Canonical $\bR/a$ & Count & Point group & Independent $A_{jl}$ \\
\midrule
Face   & $(1,0,0)$ & 6  & $C_{4v}$ & $A_{11},\; A_{22}=A_{33}$ \\
Edge   & $(1,1,0)$ & 12 & $C_{2v}$ & $A_{11}=A_{22},\; A_{33},\; A_{12}\neq 0$ \\
Corner & $(1,1,1)$ & 8  & $C_{3v}$ & $A_{11}=A_{22}=A_{33},\; A_{12}=A_{13}=A_{23}$ \\
\bottomrule
\end{tabular}
\end{center}

For the face-adjacent case, $A_{12}=A_{13}=A_{23}=0$ by $C_{4v}$
symmetry.  The edge-adjacent case has a \emph{nonzero} off-diagonal
$A_{12}$ (the two shifted axes couple).  The corner-adjacent case has
full $C_{3v}$ symmetry, reducing to only 2~independent values for
$A_{jl}$ and correspondingly few for~$B_{ijkl}$.

The delta-function collapse structure differs for each type:
\begin{itemize}[nosep]
  \item \textbf{Face}: one tent shifted by~$a$, two centred.
    Standard 2D masters $I^A_2(p,q,c)$ at $c\in\{0,1,2\}$.
  \item \textbf{Edge}: two tents shifted, one centred.
    Needs \emph{shifted} 2D masters
    $I^A_{2,\mathrm{shift}}(p,q,c,s)$ with $\sqrt{c^2+(u+s)^2+v^2}$.
  \item \textbf{Corner}: all three tents shifted.
    Double-shifted 2D masters.  But $C_{3v}$ symmetry reduces the
    count drastically.
\end{itemize}

All shifted variants follow the same sequential integration pattern
and produce exact closed forms.  The shifted master integrals for
edge and corner are already partially computed
(\texttt{CubeInterVoxelValues.wl}), but have not yet been assembled
into~$A_{jl}$,~$B_{ijkl}$,~$P_{ijkl}$.


%======================================================================
\section{Implementation plan}
\label{sec:plan}
%======================================================================

\subsection{Phase 1: Static propagators for all neighbour types}

\begin{enumerate}[label=\textbf{1\alph*.},nosep]
  \item \textbf{Face-adjacent} --- \emph{Complete.}
    \texttt{InterVoxelPropagator.wl}.

  \item \textbf{Edge-adjacent} ---
    \texttt{InterVoxelPropagatorEdge.wl} (new).
    Follow face template with two shifted tents.
    $C_{2v}$ symmetry; off-diagonal $A_{12}\neq 0$.
    New shifted 2D masters $I^A_{2,\mathrm{shift}}$.

  \item \textbf{Corner-adjacent} ---
    \texttt{InterVoxelPropagatorCorner.wl} (new).
    Three shifted tents, $C_{3v}$ symmetry.
    Double-shifted 2D masters, but heavily constrained
    (only 2~independent $A_{jl}$ values).

  \item \textbf{Export and validate.}
    Extend \texttt{CubeInterVoxelValues.wl/.json} with
    $A_{jl}$, $B_{ijkl}$, $P^{(0)}_{ijkl}$ at 16+ digit precision.
    Validate: Laplacian identity, point-group symmetry, FD cross-check.
\end{enumerate}

\subsection{Phase 2: Dynamic corrections (analytical)}

\begin{enumerate}[label=\textbf{2\alph*.},nosep]
  \item \textbf{Higher-order potentials.}
    Compute $\Psi_{00}$, $\mathrm{X}_{00}$, $\Omega_{00}$ via
    delta-function collapse for all three neighbour types.
    The kernels $|\bw|$, $|\bw|^3$, $|\bw|^5$ are smooth---these
    integrals are easier than the Newton potential.

  \item \textbf{Dynamic coefficient assembly.}
    For each order~$n$, combine the angular tensor coefficients
    ($c_p$, $c_s$ dependence) with potential derivatives to obtain
    $P^{(n)}_{ijkl}$ as exact constants.

  \item \textbf{Convergence validation.}
    Compare truncated series
    $P^{(0)} + \omega^2 P^{(1)} + \omega^4 P^{(2)}$ against WS3
    NIntegrate reference at $ka = 0.05, 0.1, 0.3, 0.5$.
\end{enumerate}

\subsection{Phase 3: Python integration}

Python's role is minimal: consume the precomputed analytical constants.
\begin{enumerate}[label=\textbf{3\alph*.},nosep]
  \item Store $P^{(n)}_{ijkl}$ as module-level constants.
  \item Neighbour classification: $\bR\to$ face/edge/corner,
    $O_h$~rotation to canonical direction, evaluate truncated series.
  \item Plug into \texttt{slab\_scattering.py} Foldy--Lax kernel
    (backward-compatible flag).
\end{enumerate}

\subsection{Phase 4: Testing}

\begin{enumerate}[label=\textbf{4\alph*.},nosep]
  \item Laplacian identity, point-group symmetry, Mathematica cross-check.
  \item $P^{\mathrm{dyn}}\to 0$ as $\omega\to 0$.
  \item Volume-averaged $\to$ point-scatterer as $|\bR|\to\infty$.
  \item Series convergence: each successive term smaller for $ka < 1$.
  \item Two-cube and slab integration tests.
\end{enumerate}


%======================================================================
\section{Summary}
\label{sec:summary}
%======================================================================

The inter-voxel strain propagator $P_{ijkl}(\bR,\omega)$ is computed
\textbf{entirely analytically} in Mathematica:
\begin{enumerate}[nosep]
  \item The UV-divergent $\order(1)$ piece is absorbed into
    derivatives of convergent potentials (Newton $\Phi_{00}$ and
    biharmonic~$\Psi_{00}$).
  \item Delta-function collapse of tent derivatives reduces 3D
    integrals to 2D masters (sequential integration $\to$ exact
    closed forms) or 1D elementary integrals.
  \item Dynamic corrections at each order $\omega^{2n}$ follow the
    same pattern with progressively \emph{smoother} kernels
    ($|\bw|^{2n-1}$), producing a convergent power series with
    analytically-computed coefficients.
\end{enumerate}
No numerical quadrature is needed at any stage.  Python only consumes
the precomputed constants for the Foldy--Lax solver.

\end{document}
